\documentclass[11pt,a4paper]{article}
\usepackage[margin=1in]{geometry}
\usepackage{enumitem}
\usepackage{booktabs}
\usepackage{hyperref}
\usepackage{xcolor}

\title{Chronological Profile: The Hidden Discovery Phase of Obelisks (2022--2023)}
\author{}
\date{}

\begin{document}
\maketitle

\section*{Overview}
Before their public announcement in early 2024, the bioinformatics discovery, algorithmic pipeline development, and laboratory validation that defined Obelisks were entirely completed during a window of intense, private research throughout late 2022 and 2023. This timeline tracks how the research team at Stanford University systematically uncovered this entirely new class of circular RNA elements.

\section*{Detailed Timeline}

\begin{enumerate}[leftmargin=*, itemsep=1.2em]
    \item \textbf{Late 2022: The Conceptual Spark \& Algorithm Design}
    \begin{itemize}[itemsep=0.3em]
        \item \textbf{Shifting the Metagenomic Paradigm:} Most historical global microbiome sequencing focused heavily on DNA. The Stanford team recognized that uncultivated, non-coding, and highly structured RNA elements represented a major blind spot in existing public datasets.
        \item \textbf{Developing ``VNom'':} To locate these hidden elements, a proprietary bioinformatics pipeline named \textit{VNom} was engineered. The tool was optimized specifically to scan massive, strand-specific RNA sequencing (RNA-seq) datasets for circularity and high self-complementarity.
    \end{itemize}

    \item \textbf{Early to Mid-2023: Mining the Human Microbiome \& Structural Naming}
    \begin{itemize}[itemsep=0.3em]
        \item \textbf{The Initial Hit (iHMP Datasets):} The pipeline was deployed against longitudinal data from the Integrative Human Microbiome Project. The team hit their first major breakthrough analyzing human stool samples, identifying a consistent $\sim$1,000-nucleotide loop that formed a rigid, symmetrical rod shape under secondary structure folding algorithms.
        \item \textbf{Coining the Term ``Obelisk'':} Because the secondary structure of these RNA loops folded into nearly perfect, self-paired rods, the team coined the name \textit{Obelisks} to mirror their monument-like physical architecture.
        \item \textbf{Discovering the Oblins:} Mid-2023 molecular modeling revealed that unlike plant viroids, which do not code for proteins, Obelisks contained functional Open Reading Frames (ORFs). They actively directed host cells to produce a completely unclassified family of proteins, which the team named \textit{Oblins}.
    \end{itemize}

    \item \textbf{Late 2023: Global Scale-Up and Host Identification}
    \begin{itemize}[itemsep=0.3em]
        \item \textbf{Massive Data Mining:} The team broadened their search parameters, utilizing a fast $k$-mer search tool called PebbleScout to mine over 3.2 million public metagenomic datasets. They successfully uncovered an astounding catalog of 29,959 distinct Obelisk varieties.
        \item \textbf{Mapping the Mouth Niche:} Data synthesis revealed that Obelisks were widespread inhabitants of the human body, present in roughly 7\% of human stool samples and a staggering 50\% of all oral microbiome profiles.
        \item \textbf{Pinpointing \textit{Streptococcus sanguinis}:} Before going public, the team matched specific donor data to establish that an oral commensal bacterium, \textit{Streptococcus sanguinis} (a primary component of early dental plaque), serves as a major biological host. Laboratory cultivation verified that the Obelisks could persist inside \textit{S. sanguinis} without disrupting native growth kinetics.
    \end{itemize}

    \item \textbf{December 2023: Manuscript Freezing}
    \begin{itemize}[itemsep=0.3em]
        \item \textbf{Finalizing the Blueprint:} By the end of 2023, the data was frozen, figures were rendered, and the comprehensive manuscript was completed. This finalized draft set up the formal preprint release on January 21, 2024, which ultimately culminated in peer-reviewed publication in the journal \textit{Cell}.
    \end{itemize}
\end{enumerate}

\section*{Relevance to Biofilm Protocols}
This timeline highlights that the primary biological reservoir for oral Obelisks is the early plaque colonizer \textit{Streptococcus sanguinis}. Because these elements are sheltered deep within the bacterial cell wall, the timeline validates mechanical biofilm disruption as the most logical pathway for managing their presence in the oral cavity.

\end{document}
